\documentclass[11pt,letterpaper]{article}
\usepackage[margin=0.92in]{geometry}
\usepackage[T1]{fontenc}
\usepackage{lmodern,microtype}
\usepackage{amsmath,amssymb,mathtools}
\usepackage[hidelinks]{hyperref}
\usepackage{enumitem}
\usepackage{xcolor}
\definecolor{graytext}{HTML}{555555}
\setlength{\parindent}{0pt}
\setlength{\parskip}{0.54em}
\setlist[itemize]{leftmargin=1.3em,itemsep=0.34em,topsep=0.25em}
\newcommand{\R}{\mathbb R}
\newcommand{\tr}{\operatorname{tr}}
\newcommand{\m}{\mathrm}
\begin{document}
\begin{center}
\url{https://github.com/xp771005/SpatialPCA-Python-Original-R}
\end{center}
\vspace{-0.45em}

\section*{1. Issues I noticed, and why}

I started with the published equations and checked their dimensions against the definitions in the text. The main issues I found seem to be in the notation for prediction, rather than in the underlying fitted model.

\textbf{Equation (3): predicting factors at new locations.} The paper describes the cross-covariance matrix $\widetilde\Sigma_l$ as $\widetilde n\times n$ (new locations by measured locations), but the displayed prediction uses $\widetilde\Sigma_l^{\top}\Sigma_l^{-1}z_l$. With $\Sigma_l^{-1}\in\R^{n\times n}$ and $z_l\in\R^n$, this multiplication generally does not conform: the first two factors have shapes $(n\times\widetilde n)(n\times n)$. Under the paper's stated convention, the transpose should be absent. The nearby cross-covariance definition also writes $K(\widetilde s,\widetilde s)$, although its two arguments should be the \emph{new} and \emph{measured} locations, $K(\widetilde s,s)$.

\textbf{Equation (4): reconstructing expression.} The displayed product uses a full loading matrix and a single predicted factor $\widetilde Z_l$. For $d>1$, a loading matrix of shape $m\times d$ cannot multiply a single $1\times\widetilde n$ factor row to produce an $m\times\widetilde n$ expression matrix. I would use the matrix $\widetilde Z$ containing \emph{all} $d$ predicted factors. Whether the covariate contribution should be added back depends on whether the target is centered expression or expression on the original scale.

\textbf{Dissertation Appendix A.1: a determinant factor.} In the dissertation's calculation that integrates out $B$, a displayed factor appears as $|X^{\top}X|^{-1/2}$. If it refers to integrating the \emph{entire} $q\times m$ matrix $B$, the factor is $|X^{\top}X|^{-m/2}$: each of the $m$ columns contributes one $|X^{\top}X|^{-1/2}$. This is a discrepancy in the dissertation expression, not yet a confirmed error in the journal's typeset supplement. With fixed $X$, it is constant in the parameters being optimized.

I also rechecked a concern I had about a missing $\log S$ in the R optimization. That concern was incorrect: the pinned original R code \emph{does} include the logarithm. I would not report it as an error.

\section*{2. How I derived and checked these results}

\textbf{New-location prediction (Eq. 3).} For one spatial factor, I wrote the joint Gaussian distribution of the measured values $z_l$ and the unmeasured values $\widetilde z_l$ as
\begin{equation}
\begin{bmatrix} z_l\\ \widetilde z_l \end{bmatrix}
\sim\mathcal N\!\left(0,
\begin{bmatrix}
\Sigma_l & \Sigma_{o*}\\
\Sigma_{*o} & \Sigma_{**}
\end{bmatrix}\right),
\qquad \Sigma_{*o}=\sigma_0^2\tau K(\widetilde s,s).
\end{equation}
The usual Gaussian conditional-mean formula gives
\begin{equation}
\boxed{\mathbb E[\widetilde z_l\mid z_l]=\Sigma_{*o}\Sigma_l^{-1}z_l},
\qquad
(\widetilde n\!\times\! n)(n\!\times\!n)(n\!\times\!1)
=\widetilde n\!\times\!1.
\end{equation}
This both gives the correct orientation and explains why the cross-kernel has to use two different sets of coordinates.

\textbf{Expression reconstruction (Eq. 4).} I went back to the original model,
\begin{equation}
Y=(XB)^\top+WZ+E.
\end{equation}
For $\widetilde n$ new positions, stacking all predicted factor rows gives $\widetilde Z\in\R^{d\times\widetilde n}$. Therefore the low-rank expression component is
\begin{equation}
\boxed{\widehat{\widetilde Y}_{\m{latent}}=\widehat W\widehat{\widetilde Z}}.
\end{equation}
If prediction is intended on the uncentered model scale, the conditional mean instead includes $(\widetilde X\widehat B)^\top$.

\textbf{Integrating out $B$ (Appendix A.1).} I defined $Y^*=(Y-WZ)^\top$, $H=X^\top X$ and $M=I_n-XH^{-1}X^\top$. Assuming $X$ has full column rank, the least-squares estimate is $\widehat B=H^{-1}X^\top Y^*$. Completing the square gives
\begin{equation}
\|Y^*-XB\|_F^2=
\|MY^*\|_F^2+
\tr\!\big[(B-\widehat B)^\top H(B-\widehat B)\big].
\end{equation}
Because the columns of $B$ separate in the second term, the Gaussian integral is
\begin{equation}
\int_{\R^{q\times m}}
\exp\!\left[-\frac{\tr((B-\widehat B)^\top H(B-\widehat B))}{2\sigma_0^2}\right]dB
=\underbrace{\big[(2\pi\sigma_0^2)^{q/2}|H|^{-1/2}\big]^m}_{(2\pi\sigma_0^2)^{qm/2}|H|^{-m/2}}.
\end{equation}
This is where the exponent $-m/2$ comes from. It assumes integration over all $B$ with a flat measure; it is not a normalized Bayesian marginal without specifying a proper prior.

Finally, I implemented the model core in Python and compared it with the pinned original R implementation on shared synthetic inputs. The recorded GitHub Actions run passed 18 tests and all nine R--Python comparisons (maximum relative error in $\tau$: $2.35\times10^{-7}$; in $WZ$: $3.17\times10^{-8}$). This supports numerical agreement for the \emph{intercept-only, full-kernel model core}; it does not test the fast approximation or real-data preprocessing. The source, results and run log are in the linked repository.

\section*{3. Full derivation, step by step}

This is the longer calculation I used to connect the model in Eq.~(1) to the parameter estimates in the implementation. I keep the number of covariates $q$ general here, even though the paper uses only an intercept ($q=1$) in its experiments. The derivation assumes $\operatorname{rank}(X)=q<n$, $K\succ0$, $\tau>0$, $W^\top W=I_d$, and $d\le m$.

\textbf{Step 1: Start from the likelihood.}
The dimensions are $Y\in\R^{m\times n}$, $X\in\R^{n\times q}$, $B\in\R^{q\times m}$, $W\in\R^{m\times d}$, and $Z\in\R^{d\times n}$. With independent errors $E_{ji}\sim\mathcal N(0,\sigma_0^2)$,
\begin{align}
Y&=(XB)^\top+WZ+E,\\
p(Y\mid B,W,Z,\sigma_0^2)
&=(2\pi\sigma_0^2)^{-mn/2}
\exp\!\left\{-\frac{\|Y-(XB)^\top-WZ\|_F^2}{2\sigma_0^2}\right\}.
\label{eq:full-likelihood}
\end{align}
Write $Y^*=(Y-WZ)^\top$. Then the squared residual becomes $\|Y^*-XB\|_F^2$.

\textbf{Step 2: Remove the covariate coefficients $B$.}
Let
\begin{equation}
H=X^\top X,\quad \widehat B=H^{-1}X^\top Y^*,\quad
M=I_n-XH^{-1}X^\top.\label{eq:full-M}
\end{equation}
The estimate $\widehat B$ follows by differentiating the residual sum of squares:
\begin{equation}
\frac{\partial}{\partial B}\|Y^*-XB\|_F^2
=-2X^\top Y^*+2X^\top XB=0.
\end{equation}
Since $M^2=M$, $M^\top=M$, and $MX=0$,
\begin{align}
Y^*-XB&=MY^*+X(\widehat B-B),\\
\|Y^*-XB\|_F^2
&=\|MY^*\|_F^2+
\operatorname{tr}\!\left[(B-\widehat B)^\top H(B-\widehat B)\right].\label{eq:full-Bsquare}
\end{align}
The cross term disappears because $X^\top M=0$. For each column $b_j$ of $B$, the Gaussian integral is
\begin{equation}
\int_{\R^q}\exp\!\left[-\frac{(b_j-\widehat b_j)^\top H(b_j-\widehat b_j)}{2\sigma_0^2}\right]db_j
=(2\pi\sigma_0^2)^{q/2}|H|^{-1/2}.
\end{equation}
There are $m$ columns, so integrating Eq.~\eqref{eq:full-likelihood} over all of $B$ gives
\begin{align}
L_B(Y\mid W,Z,\sigma_0^2)
&=\int p(Y\mid B,W,Z,\sigma_0^2)\,dB\\
&=\boxed{(2\pi\sigma_0^2)^{-m(n-q)/2}|H|^{-m/2}
\exp\!\left[-\frac{\operatorname{tr}\{(Y-WZ)M(Y-WZ)^\top\}}{2\sigma_0^2}\right].}
\label{eq:full-Bintegral}
\end{align}
I have used a flat measure for $B$ in this step. The result is an integrated-likelihood kernel, not a normalized Bayesian marginal from a proper prior. The $|H|^{-m/2}$ factor is fixed when $X$ is fixed.

\textbf{Step 3: Specify the spatial prior for $Z$.}
The rows $z_l=Z_{l\cdot}^\top\in\R^n$ are independent, with
\begin{equation}
z_l\sim\mathcal N_n(0,\sigma_0^2\tau K),\qquad l=1,\ldots,d.
\end{equation}
Their joint density is
\begin{equation}
p(Z\mid\sigma_0^2,\tau)
=(2\pi\sigma_0^2)^{-nd/2}\tau^{-nd/2}|K|^{-d/2}
\exp\!\left[-\frac{\operatorname{tr}(ZK^{-1}Z^\top)}{2\sigma_0^2\tau}\right].
\label{eq:full-prior}
\end{equation}

\textbf{Step 4: Integrate out $Z$ by completing another square.}
Multiplying Eqs.~\eqref{eq:full-Bintegral} and~\eqref{eq:full-prior} gives a $Z$-dependent quadratic
\begin{align}
Q(Z)&=\operatorname{tr}\{(Y-WZ)M(Y-WZ)^\top\}
+\tau^{-1}\operatorname{tr}(ZK^{-1}Z^\top)\\
&=\operatorname{tr}(YMY^\top)-2\operatorname{tr}(W^\top YMZ^\top)
+\operatorname{tr}\{Z(M+\tau^{-1}K^{-1})Z^\top\}.
\label{eq:full-Q}
\end{align}
In the second line I used $W^\top W=I_d$. Define
\begin{equation}
A_\tau=M+\tau^{-1}K^{-1},\qquad F_\tau=A_\tau^{-1},\qquad
C=W^\top YM.
\end{equation}
As $A_\tau$ is positive definite, the square can be completed exactly:
\begin{align}
Q(Z)
&=\operatorname{tr}\!\left[(Z-CF_\tau)A_\tau(Z-CF_\tau)^\top\right]+S(W,\tau),\\
S(W,\tau)
&=\operatorname{tr}(YMY^\top)
-\operatorname{tr}(W^\top YMF_\tau MY^\top W).
\label{eq:full-S}
\end{align}
Each of the $d$ rows contributes an $n$-dimensional Gaussian integral:
\begin{equation}
\int_{\R^{d\times n}}\exp\!\left[-\frac{Q(Z)}{2\sigma_0^2}\right]dZ
=(2\pi\sigma_0^2)^{nd/2}|A_\tau|^{-d/2}
\exp\!\left[-\frac{S(W,\tau)}{2\sigma_0^2}\right].
\end{equation}
The determinant from this integral simplifies because
\begin{align}
|A_\tau|
&=|\tau^{-1}K^{-1}|\,|I_n+\tau KM|\\
&=\tau^{-n}|K|^{-1}|I_n+\tau MK|,
\end{align}
where the last equality uses $|I+AB|=|I+BA|$. After collecting the normalizing factors, the integrated likelihood is, up to constants independent of $W$, $\tau$, and $\sigma_0^2$,
\begin{equation}
\boxed{\mathcal L(W,\tau,\sigma_0^2)\ \propto\
(\sigma_0^2)^{-m(n-q)/2}|I_n+\tau MK|^{-d/2}
\exp\!\left[-\frac{S(W,\tau)}{2\sigma_0^2}\right].}
\label{eq:full-marginal}
\end{equation}

\textbf{Step 5: Estimate $\sigma_0^2$, $W$, and $\tau$.}
For fixed $W$ and $\tau$, put $N=m(n-q)$ and $v=\sigma_0^2$. The log-likelihood terms involving $v$ are
\begin{equation}
\ell(v)=-\frac N2\log v-\frac{S(W,\tau)}{2v}+\mathrm{const}.
\end{equation}
Differentiating and setting the result to zero,
\begin{equation}
-\frac{N}{2v}+\frac{S(W,\tau)}{2v^2}=0
\qquad\Longrightarrow\qquad
\boxed{\widehat\sigma_0^2=\frac{S(W,\tau)}{m(n-q)}}.
\label{eq:full-sigma}
\end{equation}
For fixed $\tau$, define the symmetric, positive-semidefinite matrix
\begin{equation}
G_\tau=YMF_\tau MY^\top\in\R^{m\times m}.
\end{equation}
Minimizing $S$ is equivalent to maximizing $\operatorname{tr}(W^\top G_\tau W)$ subject to $W^\top W=I_d$. By the Rayleigh--Ritz/Ky Fan principle, the columns of $\widehat W_\tau$ span the top $d$ eigenspace of $G_\tau$. Thus
\begin{equation}
S_*(\tau)=\operatorname{tr}(YMY^\top)-\sum_{l=1}^d\lambda_l(G_\tau),
\end{equation}
where the eigenvalues are sorted from largest to smallest. Substituting $\widehat\sigma_0^2$ into Eq.~\eqref{eq:full-marginal} gives the one-dimensional objective
\begin{equation}
\boxed{J(\tau)=\frac{m(n-q)}2\log S_*(\tau)
+\frac d2\log|I_n+\tau MK|,\qquad
\widehat\tau=\arg\min_{\tau>0}J(\tau).}
\label{eq:full-profile}
\end{equation}
In practice, $\eta=\log\tau$ is convenient for numerical optimization. This derivation also explains why a $\log S_*$ term belongs in the R/Python objective; the original R implementation indeed contains it.

\textbf{Step 6: Recover the spatial factors and covariate coefficients.}
From the completed square in Eq.~\eqref{eq:full-Q}, the posterior distribution for each row (after integrating $B$) is
\begin{equation}
z_l\mid Y,W,\tau,\sigma_0^2\sim
\mathcal N_n\!\left(F_\tau MY^\top w_l,\,\sigma_0^2F_\tau\right).
\end{equation}
Therefore the posterior-mean spatial PCs and the estimated covariate coefficients are
\begin{align}
\boxed{\widehat Z=\widehat W^\top YM F_{\widehat\tau}},\qquad
\boxed{\widehat B=(X^\top X)^{-1}X^\top
(Y-\widehat W\widehat Z)^\top}.
\label{eq:full-ZB}
\end{align}
The individual columns of $W$ need not match between two runs because an orthogonal rotation of the factors leaves $WZ$ unchanged. That is why I compared the loading subspace and reconstructed $WZ$ in the R/Python tests.

\textbf{Step 7: Predict at unobserved locations.}
Let $K_{oo}=K(s,s)\in\R^{n\times n}$ and $K_{*o}=K(\widetilde s,s)\in\R^{\widetilde n\times n}$. The joint prior for a factor at measured and new locations is Gaussian. Conditioning on the measured factor values gives
\begin{align}
\mathbb E[\widetilde z_l\mid z_l]
&=(\sigma_0^2\tau K_{*o})(\sigma_0^2\tau K_{oo})^{-1}z_l\\
&=\boxed{K_{*o}K_{oo}^{-1}z_l}.
\end{align}
Applying this map to all $d$ estimated factor rows produces $\widehat{\widetilde Z}\in\R^{d\times\widetilde n}$. The low-rank predicted expression is $\widehat W\widehat{\widetilde Z}$; if we want the model's uncentered conditional mean, we also add $(\widetilde X\widehat B)^\top$. This is why I interpreted Eqs.~(3)--(4) as notation issues in Section 1 rather than evidence of incorrect fitted parameters.

\textbf{Step 8: A note on the computational form.}
The determinant and inverse can be evaluated without repeatedly inverting $K^{-1}$. Write $P=I_n+\tau K$ and $A_0=I_n+\tau^{-1}K^{-1}$; then
\begin{align}
|I_n+\tau MK|
&=|P|\,\frac{|X^\top P^{-1}X|}{|X^\top X|},\label{eq:full-det}\\[2pt]
F_\tau
&=A_0^{-1}+A_0^{-1}X
\big(X^\top X-X^\top A_0^{-1}X\big)^{-1}
X^\top A_0^{-1}.\label{eq:full-woodbury}
\end{align}
The second identity follows from Woodbury after writing $M=I_n-X(X^\top X)^{-1}X^\top$. If $K=UDU^\top$ with $U$ square and orthogonal, the simpler inverse is
\begin{equation}
A_0^{-1}=U\,\operatorname{diag}\!\left(
\frac{\tau\delta_i}{1+\tau\delta_i}\right)U^\top.
\end{equation}
One should not substitute the inverse of a truncated rank-$r$ kernel $K_r=U_rD_rU_r^\top$ into these full-rank identities: $K_r$ is singular when $r<n$. A model using a rank-$r$ Gaussian prior instead leads to the restricted operator
\begin{equation}
F_r=U_r\big(\tau^{-1}D_r^{-1}+U_r^\top MU_r\big)^{-1}U_r^\top.
\end{equation}
I treat this as a separate approximation, not as an error in the original software. The GitHub comparison reported above tested the \emph{full-kernel} case.

\end{document}
